\begin{frame}{How is data generated}
    \centering
    \begin{figure}
        \centering
        \includegraphics[width=\textwidth,height=0.45\textheight,keepaspectratio]{img/eiger2.jpg}
        \caption{Eiger Detector}
        \label{fig:my_label}
    \end{figure}
\end{frame}
\note{Diamond can be viewed as  giant video camera  filming a sample and it rotates on its axis. We do this using Fancy camera's known as detectors, now these camera's generate a incredible amount of unstructured data a second. Put to put the amount into perceptive, a Netflix movie your might stream for a evenings entertainment is about 7GB whereas our detectors are producing anywhere between 1GB/s to 40GB/s.}
\begin{frame}{Diamond = Netflix? }
\begin{columns}
    \begin{column}{0.30\textwidth}
    \begin{figure}
        \centering
        \includegraphics[width=\textwidth,height=0.55\textheight,keepaspectratio]{img/Diamond.eps}
        \caption{Diamond}
        \label{fig:my_label}
    \end{figure}
    \end{column}
    \begin{column}{0.30\textwidth}
    \begin{figure}
        \centering
        \includegraphics[width=\textwidth,height=0.55\textheight,keepaspectratio]{img/equals (2).png}
        \label{fig:my_label}
    \end{figure}
    \end{column}
    \begin{column}{0.30\textwidth}
    \begin{figure}
        \centering
        \includegraphics[width=\textwidth,height=0.55\textheight,keepaspectratio]{img/netflix.jpg}
        \caption{Netflix}
        \label{fig:my_label}
    \end{figure}
    \end{column}
\end{columns}
\end{frame}
\note{At Diamond we stream science instead of movies}